\chapter{REFLEXIÓN CRÍTICA PERSONAL}

El caso Amazon Recruiting Tool constituye, a mi juicio, un ejemplo especialmente valioso para comprender que la discriminación algorítmica rara vez es intencional: surge, con mayor frecuencia, de la combinación entre datos históricos sesgados, decisiones técnicas aparentemente neutrales y una confianza excesiva en la supuesta objetividad de los sistemas de inteligencia artificial. El principal error de Amazon no fue, en mi opinión, haber intentado automatizar la selección de personal —un objetivo legítimo desde la perspectiva de la eficiencia—, sino haber subestimado la necesidad de auditar de forma sistemática y continua los resultados del modelo por subgrupos demográficos antes y durante su uso, así como no haber diversificado suficientemente tanto los datos de entrenamiento como el propio equipo de desarrollo.

Entre las medidas que debieron haberse adoptado destacan: la realización de una evaluación de impacto ético y de sesgos antes del despliegue del sistema, incluso en ausencia de una obligación legal explícita en aquel momento; la incorporación de técnicas de mitigación de sesgos para neutralizar variables proxy de género; la creación de mecanismos de auditoría periódica por parte de un comité interno de ética de datos; y una mayor transparencia interna que hubiera permitido detectar el problema antes de que transcurriera un año completo de funcionamiento sesgado.

La principal enseñanza que este caso deja de cara al futuro es que la neutralidad algorítmica es, en gran medida, un mito si no se somete a verificación activa: un sistema puede parecer objetivo precisamente porque oculta, bajo una fachada de cálculo matemático, los mismos sesgos que ya existían en la sociedad y en los datos que la reflejan. Además, el caso confirma que la responsabilidad no puede delegarse en la máquina: son las organizaciones, y no los algoritmos, quienes deben responder por las decisiones automatizadas que despliegan, tal como reconocen tanto el AI Act como las guías más recientes de la EEOC \cite{eeoc2023}.
